\section{Introduction}
\label{sec:introduction}

Electricity price forecasting occupies an unusual position in forecasting research. Forecast errors are statistically measurable, but the practical value of reducing them depends on the decisions made with the forecasts. This distinction is especially important in day-ahead electricity markets, where prices can be negative, exhibit abrupt spikes, and respond nonlinearly to load, renewable generation and market conditions. These features have motivated a large literature spanning autoregressive models, regularized regression, machine learning and probabilistic forecasting \citep{weron2014,zielweron2018,lago2021}. Yet improvement in a conventional accuracy measure such as mean absolute error (MAE) is not, by itself, evidence that a forecast produces a better economic decision.

The distinction between statistical and economic forecast quality is established prior work. \citet{kathziel2018} explicitly connect electricity-price forecast quality to realized economic gains, while \citet{serafin2025} show that choices affecting forecast loss can materially alter both trading profit and risk. Most recently, \citet{maciejowska2026} compare 192 forecasting specifications in the German/Germany--Luxembourg market and show that RMSE and MAE are only weakly related to battery-arbitrage income. The same issue extends to probabilistic forecasts: \citet{nitka2023} find that a statistically better predictive-distribution combination need not generate higher realized profit. The relevant research question is therefore not whether accuracy and economic value can diverge, but whether a forecasting advantage remains useful when the information set, uncertainty method and economic decision rule are all constrained ex ante.

That requirement is more demanding than it appears. A day-ahead forecasting exercise is economically meaningful only if each predictor could have been available when the simulated decision was made. Historical electricity databases often contain load, wind, solar and other variables that are highly informative ex post. Their presence in a historical API, however, does not establish that the particular stored value was publicly available at the relevant forecast origin. Publication deadlines, revisions and data vintages therefore matter. Prior work supports the predictive usefulness of fundamentals and weather information \citep{kathziel2018,sgarlato2023}, but that evidence does not prove that every historical ENTSO-E renewable forecast used in a new application was available at a specific pre-auction timestamp.

A second difficulty concerns evaluation design. Electricity-price models are typically developed through repeated choices over features, algorithms, hyperparameters, calibration windows and error measures. Unless these choices are separated from the final evaluation period, reported out-of-sample performance can become contaminated by iterative model development. The best-practice literature therefore emphasizes chronological testing, strong benchmarks, common evaluation samples and model selection that does not use the final test period \citep{weron2014,lago2021}. This discipline is even more important when forecasts feed an economic strategy, because thresholds, costs, uncertainty filters and operating rules can themselves be tuned after observing profits.

This study addresses these issues through an end-to-end forecast-to-decision experiment for the Germany--Luxembourg (DE-LU) day-ahead electricity market. The empirical sample starts on 1 January 2019, after the German--Austrian bidding-zone split. The simulated forecast and decision origin is fixed at 11:45 Europe/Berlin on D-1. Public ENTSO-E price and fundamental data are cleaned under explicit local-delivery-day logic, including daylight-saving-time transitions and the move of Single Day-Ahead Coupling (SDAC) to 15-minute market time units for delivery from 1 October 2025 \citep{entsoe2025}.

The predictor set is divided into two information levels. The \emph{Full} specification contains lagged prices, calendar effects, the day-ahead load forecast, renewable-generation forecasts and variables derived from them. The \emph{Tier-1} specification removes the wind- and solar-derived variables whose precise historical public availability at 11:45 D-1 cannot be established from the historical ENTSO-E interface. This distinction does not assert that those variables contain look-ahead information. It identifies an evidential boundary: their exact vintage at the simulated forecast origin is not reconstructable from the data source used here.

All model development is restricted to 2019--2025. Naive 24-hour and 168-hour lags provide transparent benchmarks, ElasticNet provides a regularized linear benchmark, and XGBoost is evaluated as a nonlinear challenger. Model selection is chronological. Forecast uncertainty is estimated from out-of-sample residuals and corrected to respect the common whole-day D-1 forecast origin: every hour of delivery day \(D\) uses residual information from delivery dates strictly earlier than \(D\). Because the selected residual offsets are common to all hours within a delivery day, the resulting uncertainty-aware economic rule is algebraically a day-varying abstention threshold rather than an alternative intraday pair-selection rule.

Economic value is evaluated with a deliberately transparent storage-arbitrage rule rather than a full battery-dispatch model. Each day allows at most one charge interval and one later discharge interval. Six strategies distinguish no trading, lag-based decisions, Full and Tier-1 XGBoost decisions, and uncertainty-filtered variants. Before the holdout is opened, the primary efficiency and degradation assumptions, the entire \(3\times3\) sensitivity grid, the uncertainty procedure and the primary economic confirmation rule are frozen.

The final evaluation uses a first-look holdout from 1 January through 31 July 2026. Its structural completeness was checked before outcome exposure, but its realized prices were not used for model selection, tuning, uncertainty-window choice, strategy design or economic interpretation. The final forecasting models, uncertainty method and economic rules were frozen before the first 2026 performance metric was inspected.

The holdout produces three central findings. First, Full XGBoost retains a large predictive advantage: MAE is \EUR{17.51}/MWh compared with \EUR{29.33}/MWh for lag-24 and \EUR{24.36}/MWh for Tier-1 XGBoost. Second, the forecasting improvement translates into economic value under the pre-specified decision rule. At the primary parameter setting, the Full XGBoost point strategy earns \EUR{20,002.02} versus \EUR{18,564.42} for lag-24, and the incremental value remains positive in all nine pre-specified efficiency--cost cells. Third, uncertainty filtering is not unambiguously profit-enhancing. The Full uncertainty-aware strategy gives up \EUR{872.01} of cumulative net P\&L but increases the trade hit rate from 91.9\% to 99.4\% and improves the worst daily result from \EUR{-17.55} to \EUR{-0.83}.

The contribution is consequently methodological and evidential rather than algorithmic. The paper combines: (i) an explicit point-in-time information boundary through Full and Tier-1 specifications; (ii) a pre-specified protocol frozen before holdout exposure; (iii) a genuine first-look 2026 holdout; (iv) a pre-specified nine-cell economic robustness criterion; (v) delivery-day-safe uncertainty implemented transparently as an abstention threshold; and (vi) economic and tail-risk diagnostics. The study does not claim that XGBoost is universally superior, that lower MAE necessarily implies greater economic value, or that the stylized storage results constitute deployable investment returns.

The remainder of the paper is organized as follows. Section~\ref{sec:literature} positions the study in the forecasting, uncertainty and storage-value literature. Section~\ref{sec:data} describes the market, data and information-set construction. Section~\ref{sec:methods} presents forecasting, uncertainty and economic methods. Section~\ref{sec:development} reports development evidence and the freeze decisions. Section~\ref{sec:holdout} presents the frozen 2026 holdout. Section~\ref{sec:discussion} discusses interpretation and limitations, and Section~\ref{sec:conclusion} concludes.
